\section{Evaluation}
\label{sec:evaluation}

This section reports an artifact-based assessment of the workflow, focusing on (i) embedding capacity of the selection channel, (ii) naturalness and distributional divergence of successful stego samples, and (iii) operational reliability under the sender-side verification loop. We center results on a clean run artifact (\texttt{run\_id} \texttt{20260509T203344Z}, directory \texttt{metrics/e2e\_runs/fresh\_metrics\_200\_20260509T233342Z}) and explicitly separate generation/decode failures from infrastructure errors.

\subsection{Capacity Analysis}
The channel capacity of Layer~2 is fundamentally dependent on the size of the shared, deterministically reconstructed angle set. If the reconstructed set contains $N$ Angles, then Layer~2 contributes $\log_2 N$ bits per message instance (rounded up to an integer bit-width in the implementation). In practice, broader threads and richer research corpora tend to yield larger and more diverse angle sets than narrowly constrained topics, increasing available capacity.

\subsection{Perceptual \& Statistical Imperceptibility}
Imperceptibility is evaluated through the dual lenses of semantic coherence and contextual integration. Unlike traditional steganographic methodologies that risk producing disjointed or statistically anomalous text, our intent-based approach aims to ensure that each stego contribution maintains a clear thematic alignment with the surrounding discourse. In this work we report GPT-2 perplexity and word-unigram KL/JSD as coarse statistical proxies, noting that these metrics do not fully capture long-context semantic quality \cite{wang2022perplexity} and do not replace dedicated steganalysis benchmarks.

Table~\ref{tab:method-metrics} summarizes perplexity (GPT-2) and word-unigram divergence against two baselines: the matched primary post and a global comment corpus. Values are computed over the successful, receiver-validated stego samples produced in run \texttt{20260509T203344Z} (\texttt{requested}=200, \texttt{succeeded}=183).

\begin{table}[t]
	\centering
	\caption{Perplexity and distribution divergence for stego comments (GPT-2 PPL; KL/JSD on word unigrams).}
	\label{tab:method-metrics}
	\small
	\begin{tabularx}{\columnwidth}{@{}l l X c X@{}}
		\toprule
		Category & Metric & Baseline / context & Mean & Range (min / max) \\
		\midrule
		Perplexity & PPL (GPT-2) & 183 successful stego strings & $\approx 142.3$ & 22.8 -- 653.8 \\
		\midrule
		Divergence & KL & Matched primary post & $\approx 9.24$ & --- \\
		Divergence & KL & Global comment corpus & $\approx 6.50$ & --- \\
		Divergence & JSD & Matched primary post & $\approx 0.590$ & --- \\
		Divergence & JSD & Global comment corpus & $\approx 0.588$ & --- \\
		\bottomrule
	\end{tabularx}
\end{table}

\paragraph{Definitions.}
Perplexity is computed from average per-token log-probability,
$2^{-\frac{1}{n} \sum_{i=1}^{n} \log_2 P(x_i \mid x_{<i})}$:
lower PPL indicates more natural continuations under the model; higher PPL indicates surprising or awkward text.
The observed maximum ($\sim$654) likely reflects a small number of pathological samples that are highly unnatural to GPT-2.

Kullback--Leibler (KL) divergence compares stego word unigrams $P_S$ to a baseline $P_C$ via
$\sum_w p_{\text{stego}}(w) \log \frac{p_{\text{stego}}(w)}{p_{\text{baseline}}(w)}$;
lower KL indicates closer statistical match to the baseline.
KL being lower against the global corpus than against the matched post suggests stego wording drifts toward generic phrasing rather than thread-specific vocabulary.

Jensen--Shannon divergence (JSD) is the symmetric, bounded counterpart,
$D_{\mathrm{JS}}(P_C \,\|\, P_S) = \frac{1}{2} D_{\mathrm{KL}}\bigl(P_C \,\|\, \frac{P_C + P_S}{2}\bigr) + \frac{1}{2} D_{\mathrm{KL}}\bigl(P_S \,\|\, \frac{P_C + P_S}{2}\bigr)$,
often used to upper-bound distinguishability from the true language distribution.
Word pairs absent from one side of the comparison are handled with Laplace-style smoothing (e.g., $\alpha = 10^{-6}$) so KL and JSD remain well defined.

Table~\ref{tab:metrics-vs-literature} contrasts these means with representative values from related work (exact comparability varies by corpus and tokenization).
Table~\ref{tab:slr-metrics-comparison} uses the same column layout as our systematic-review extraction sheet: paper, model, corpus, reported metrics, and context tags, with our row listed first for side-by-side reading.

\begin{landscape}
	\footnotesize
	\setlength{\tabcolsep}{4pt}
	\begin{longtable}{@{}p{2.05cm}p{2.35cm}cp{1.85cm}p{5.35cm}p{1.45cm}p{1.55cm}p{1.95cm}@{}}
		\caption{This work versus prior papers (SLR-style summary; metrics are as stated in each paper and may use different models, tokenizations, and baselines).}
		\label{tab:slr-metrics-comparison} \\
		\toprule
		Paper & LLM & Yr. & Dataset & Result & Ctx.\ aware & Cat.\ ctx.\ & Repr.\ ctx.\ \\
		\midrule
		\endfirsthead
		\multicolumn{8}{c}{\tablename\ \thetable{}---continued from previous page} \\
		\toprule
		Paper & LLM & Yr. & Dataset & Result & Ctx.\ aware & Cat.\ ctx.\ & Repr.\ ctx.\ \\
		\midrule
		\endhead
		\midrule
		\multicolumn{8}{r}{\emph{Continued on next page}} \\
		\endfoot
		\bottomrule
		\endlastfoot
		This work & Configurable black-box LLMs (deterministic decode; stochastic synthesis); GPT-2 for PPL & 2026 & Threaded forum posts; run \texttt{20260509T203344Z} (183 successes) & PPL $\approx 142.3$; KL 6.50 (global) / 9.24 (matched); JSD 0.588 (global) / 0.590 (matched) & explicit & social thread & selection channel (reply + angle) \\
		\midrule
		VAE-Stega \cite{yang2020vae} & BERT--LSTM; LSTM--LSTM (from scratch) & 2020 & Twitter (2.6M sents.); IMDB (1.2M sents.) & PPL 28.879; $\Delta$MP 0.242; KLD 3.302; JSD 10.411; Acc 0.600 & non-explicit & pre-text & text \\
		General MLM RH \cite{zheng2022general} & BERT-base & 2022 & BookCorpus & BPW 0.5335; F1 0.9402; PPL 134.22 & non-explicit & pre-text & text \\
		Co-Stega \cite{liao2024co} & Llama-2-7B/13B-chat; fine-tuned GPT-2 & 2024 & Twitter; tweet data for fine-tuning & SR1 60.87\%; SR2 98.55\%; cap.\ 44.91 bit; BPW 2.31; PPL 16.75; SimCSE 0.69 & explicit & social media & text \\
		Joint BERT--GAT \cite{ding2023joint} & LSTM+attn.; GAT; BERT MLM & 2023 & OPUS & PPL 13.917; KLD 2.904; SIM 0.812; ER 0.365; best Acc 0.575 & explicit & pre-text & text \\
		Discop \cite{ding2023discop} & GPT-2 & 2023 & IMDB & $p{=}1.00$; ave.\ KLD 0; max KLD 0; cap.\ 5.76 bit/token; PPL $\approx 42$ (stego quality in paper) & non-explicit & tuning + pretext & text \\
		LLM generative stego \cite{wu2024generative} & Any LLM (e.g., GPT-4 in eval.) & 2024 & [not specified] & len.\ 13.3 words; BPW 5.93; PPL 165.76; SS 0.5881; BiLSTM Acc 49.2\% & explicit & [varies] & [varies] \\
		Zero-shot ling.\ stego \cite{lin2024zero} & LLaMA2-Chat-7B; GPT-2 for baselines & 2024 & IMDB; Twitter & PPL 8.81; JSD$_\text{full}$ $17.9{\times}10^{-2}$; JSD$_\text{half}$ $16.86{\times}10^{-2}$; JSD$_\text{zero}$ $13.40{\times}10^{-2}$ & explicit & zero-shot + prompt & text \\
		Provably secure disamb.\ \cite{qi2024provably} & LLaMA2-7B; Baichuan2-7B & 2024 & IMDb (EN/ZH samples) & err.\ 0\%; ave.\ KLD 0; ave.\ PPL 3.19 (EN), 7.49 (ZH); cap.\ 1.03--3.05 bit/token & non-explicit & pretext & text \\
		Hi-Stega \cite{wang2023hi} & GPT-2 & 2024 & Yahoo! News (titles/bodies/comments) & PPL 109.60; MAUVE 0.2051; ER2 10.42; $\Delta$cosine 0.0088; $\Delta$SimCSE 0.0191 & explicit & social media & text \\
	\end{longtable}
\end{landscape}

\begin{table}[t]
	\centering
	\caption{Comparison of measured means to illustrative values from the literature.}
	\label{tab:metrics-vs-literature}
	\footnotesize
	\begin{tabularx}{\columnwidth}{@{}l c X@{}}
		\toprule
		Metric & Ours (mean) & Source reference and context \\
		\midrule
		PPL & 142.3 & 134.11 (human mean), VAE-Stega / IMDB training baseline. \\
		PPL & & 20.92 (RNN-Stega), VAE-Stega RNN at 1.0 bpw. \\
		PPL & & 165.76 (LLM-Stega), black-box generation with GPT-4. \\
		PPL & & 42.23 (Discop), stego quality with GPT-2. \\
		\midrule
		KL & 6.50--9.24 & 19.51 (RNN-Stega), VAE-Stega RNN on IMDB. \\
		KL & & 0.045 (Meteor), high-entropy English sampling. \\
		KL & & 0.00 (Discop), distribution-matching stego. \\
		\midrule
		JSD & 0.588--0.590 & 14.48 (RNN-Stega), VAE-Stega RNN on IMDB. \\
		JSD & & 0.18 (Zero-shot LLM), in-context learning ($17.9 \times 10^{-2}$). \\
		\bottomrule
	\end{tabularx}
\end{table}

\paragraph{Reading the comparison.}
Mean PPL ($\approx 142.3$) is of similar scale to the human-text average reported for VAE-Stega (134.11) and lower than a cited black-box LLM-Stega reference (165.76), i.e., within a plausible range for non-adversarial generated forum-style text when evaluated under GPT-2.
KL (6.50--9.24) is much smaller than the RNN-Stega IMDB figure (19.51), indicating stronger unigram-level alignment with the chosen baselines than older recurrent generators, but larger than methods that tightly match a model distribution (Meteor, Discop), which operate at token level or with provable distribution matching.
JSD (0.588--0.590) exceeds the cited zero-shot LLM value (0.18), so the stego text retains a more distinct unigram signature than that reference.
Finally, lower KL to the global corpus than to the matched post parallels a known pattern: as surface quality is pushed, text may move toward background forum language and away from idiosyncratic thread-specific vocabulary.

\subsection{System Robustness \& Recovery Rate}
The sender-side verification loop improves operational reliability by rejecting candidates that do not decode back to the selected Angle. In run \texttt{20260509T203344Z} (\texttt{requested}=200), the pipeline produced 183 successful, receiver-validated samples and 17 failures after exhausting the retry budget. The primary failure modes were (i) decoding validation failures (10/17) and (ii) malformed synthesis output that did not match the required JSON schema (7/17). These failures highlight that reliability is primarily limited by the ability to synthesize a decodable witness under stochastic generation, rather than by payload recovery after validation.

\subsection{Resistance to Automated Steganalysis}
The design objective of the framework is to reduce semantic and pragmatic anomalies by conditioning synthesis on the full thread context and by encoding information in conversational decisions (reply targeting and Angle selection) rather than local token artifacts. We therefore expect detection to require thread-level pragmatic reasoning rather than relying solely on token-level distribution cues \cite{Wang_2023}. A dedicated steganalysis benchmark evaluation is left to future work; here we limit claims to the measured distributional and operational metrics reported above.
